%%%% LISTA DE ABREVIATURAS, SIGLAS E ACRÔNIMOS
%%
%% Relação, em ordem alfabética, das abreviaturas (representação de uma palavra por meio de alguma(s) de sua(s) sílaba(s) ou
%% letra(s)), siglas (conjunto de letras iniciais dos vocábulos e/ou números que representa um determinado nome) e acrônimos
%% (conjunto de letras iniciais dos vocábulos e/ou números que representa um determinado nome, formando uma palavra pronunciável).
%%
%% Este arquivo para definição de abreviaturas, siglas e acrônimos é utilizado com a opção "file" (arquivo)

\listadeabrevsiglaseacr%% Formatação da lista de abreviaturas, siglas e acrônimos

%% ATENÇÃO - não use esse arquivo, utilize o lista-simbolos.tex

% \begin{listadeabreviaturas}%% Ambiente listadeabreviaturas (Abreviaturas: Representação & Definição \\)
% %% ATENÇÃO - não use esse arquivo, utilize o lista-simbolos.tex
% %% as abreviaturas estão lá também...
% art. & Artigo   \\
% cap. & Capítulo \\
% sec. & Seção    \\

% \end{listadeabreviaturas}

\begin{listadesiglas}%% Ambiente listadesiglas (Siglas: Representação & Definição \\)
%% ATENÇÃO - não use esse arquivo, utilize o lista-simbolos.tex
%ABNT  & Associação Brasileira de Normas Técnicas                              %     \\

AI   & \textit{Artificial Intelligence} \\
ALU  & \textit{Arithmetic Logic Unit} \\
CISC & \textit{Complex Instruction Set Computer}\\
CPI  & \textit{Cycles Per Instruction} \\
CU   & \textit{Control Unit} \\
FPGA & \textit{Field-Programmable Gate Array} \\
HDL  & \textit{Hardware Description Language} \\
HPC  & \textit{High-Performance Computing} \\
I/O  & \textit{Input/Output} \\
ISA  & \textit{Instruction Set Architecture} \\
LLM  & \textit{Large Language Models} \\
PC   & \textit{Program Counter} \\
RISC & \textit{Reduced Instruction Set Computer} \\

\end{listadesiglas}

%\vspace{1cm}\noindent\textcolor{red}{Elemento opcional. Relação em ordem alfabética das abreviaturas e siglas presentes no texto, seguidas do respectivo significado.}

%\begin{listadeacronimos}%% Ambiente listadeacronimos (Acrônimos: Representação & Definição \\)
%Gimp & Programa de Manipulação de Imagem GNU, do inglês \textit{GNU Image Manipulation Program} \\
%\end{listadeacronimos}

%% Obs.: comente ou remova as linhas correspondentes ao ambiente que se deseja remover desta lista.
